\section{Limitations and ethical considerations}
\label{sec:limitations}
\textbf{Scope.} We study one dataset and two conventional ImageNet-pretrained backbones that differ in
capacity and pretraining data (Table~\ref{tab:models}); differences between them are not attributed to
inductive bias. Faithfulness depends on the perturbation substrate and step size; we use one fixed
setting. The pointing game and IoU depend on the tolerance and threshold; we report two tolerances and
one threshold. Three seeds give a coarse estimate of training variation.
\textbf{Data.} HAM10000 is a public, de-identified dataset collected at sites in Austria and
Australia~\cite{tschandl2018ham10000}. It has no skin-type annotations, so we cannot assess performance
or explanation quality across skin tones. \textbf{Risk of misleading explanations.} A heatmap that lands on a centered lesion can look convincing to a
clinician even when a model-free prior would land in the same place; relying on such a heatmap as evidence of
model reasoning could create unwarranted trust. We state this as a risk motivated by our localization results,
not as a measured effect on clinicians; we ran no user study. \textbf{Use.} This is a research evaluation of
explanation benchmarks; nothing here supports clinical use of these models or their explanations.
